\begin{enumerate}[leftmargin=1.8em,itemsep=2pt]
  \item \textbf{Tor OpenRouter nie powstał.} Plan przewidywał 3 tory pomiaru P4.
  Zabezpieczenie komputera autora zablokowało odczyt klucza dostępu
  \pkq{[obserwacja: 2026-11-12, 15:15]}. Pomiar ma 2 tory: lokalny i Claude.
  \item \textbf{Zbiór ma \pkItems{} pozycji.} Plan przewidywał 40 pozycji. Dodatkowe
  pozycje to 2 luki standardu (G1, G2) i 2 pozycje o słowach kluczowych (K1, K2).
  \item \textbf{Kontrola kotwic.} Kontroli kotwic nie przeszło \pkProbeFailed{}
  z \pkProbeLocal{} modeli lokalnych. Wynik takiego modelu stoi w tabeli \ref{tab:probe}
  i nie wchodzi do miar zgodności.
  \item \textbf{Tor Claude nie ma kotwic.} Budżet toru wynosił co najwyżej 600 wywołań,
  więc tor obejmuje tylko sondy strefy szarej w 2 próbach.
  \item \textbf{Przerwa toru Claude.} O godz. 18:06 polecenie \texttt{claude} zniknęło na
  czas aktualizacji pakietu i bieg zakończył się błędem. Autor dodał do skryptu ponawianie
  i wznowił bieg; skrypt zachował odpowiedzi sprzed przerwy
  \pkq{[obserwacja: 2026-11-12, 18:07]}. Bieg przerwany wykonał 34 zapytania i nie ma
  wiersza w tabeli \ref{tab:runs}.
  \item \textbf{Odczyt werdyktu w torze Claude.} W 4 sondach model napisał rachunek przed
  werdyktem, a skrypt odczytywał werdykt z początku odpowiedzi. Autor zmienił regułę:
  werdyktem jest końcowe słowo odpowiedzi. Autor powtórzył te 4 sondy (8 wywołań)
  i przeliczył 1 wynik z zapisanej odpowiedzi. Tor Claude zużył razem 552 wywołania.
  \item \textbf{Tokenizery modeli Bielik.} Serwis Hugging Face zwrócił błąd 401, więc
  pomiar P3 pomija 2 tokenizery.
  \item \textbf{Pobieranie modeli.} Łącze dało prędkość 1,3 MB/s \pkq{[obserwacja:
  2026-11-12, 17:41]}. Program \texttt{ollama} zatrzymał pobieranie 1 raz; skrypt ponawia
  pobieranie co 10 minut. Modele zajmują miejsce na partycji \texttt{/mnt/windows}, która
  przed pobraniem miała 18 GB wolnego miejsca.
  \item \textbf{Stan modeli lokalnych.} Pomiar objął \pkProbeLocal{} modele lokalne z 4 planowanych: \texttt{llama3.2:3b},
\texttt{gemma3:4b}, Bielik 1.5B i Bielik 4.5B. Pobranie 4 modeli zajęło 2 h 42 min
\pkq{[obserwacja: 2026-11-12, od 17:27 do 20:09]}.

  \item \textbf{Parser bez poprawki Leo.} Gramatyka PK w pomiarze P2 ma wykładnik
  \pkParseExpPk{}, a nie 1. Czas liniowy wymaga parsera z poprawką Leo.
  \item \textbf{Sprawdzarka jest przybliżona.} Sprawdzarka zgłasza także fałszywe
  trafienia, na przykład przymiotnik w nazwie \pkq{małe litery}. Precyzja na próbie
  \pkPrecN{} trafień wynosi \pkPrec{} \%. Czułość jest nieznana.
  \item \textbf{Podział aktów prawnych na zdania.} Akt składa się z jednostek
  redakcyjnych. Skrypt usuwa numery jednostek i dzieli tekst regułą interpunkcyjną, więc
  mediana długości zdania jest przybliżeniem.
\end{enumerate}
